\documentclass[11pt,a4paper]{article}
\usepackage[T1]{fontenc}
\usepackage{lmodern,amsmath,amssymb,geometry}
\geometry{margin=29mm}
\title{EPPA--III: the complementary-half partition lemma}
\author{Independent validation supplement}
\date{9 October 2026}
\begin{document}
\maketitle

\noindent
\textbf{Purpose.} Proposition 3.23 invokes Lemma 3.7 without
its hypothesis that \(H\setminus G\cong sK_t\).
The following lemma provides that hypothesis under
precisely the non-homogeneity assumption in Proposition 3.23.
The circulation manuscript remains untouched.

\paragraph{Lemma.}
Let \(H\) be a finite non-homogeneous vertex-transitive
EPPA witness for an induced \(G=sK_t\) such that
\(|H|=2st\), with \(s,t\ge1\).
Then \(H[V(H)\setminus V(G)]\cong sK_t\).

\paragraph{Proof.}
Put \(B=V(H)\setminus V(G)\), and \(n=st=|B|\).
Since \(H\) is regular and \(|G|=|B|\),
double-counting the edges crossing \((G,B)\) gives
\[
                 |E(H[G])|=|E(H[B])|.        \tag{*}
\]
If \(s=1\), this is the maximum possible number
\(\binom t2\) of edges on \(B\), so \(H[B]=K_t\).
If \(t=1\), both sides have zero edges, so \(H[B]=sK_1\).
Assume \(s,t\ge2\).

Let \(W=\operatorname{Aut}(G)=S_t\wr S_s\).
For \(w\in B\), let
\[
             P(w)=(N_H(w)\cap C_i)_{i=1}^s.
\]
Every member of \(W\) extends, by EPPA, to an
automorphism of \(H\) preserving both \(G\) and \(B\).
Consequently the \(W\)-orbit of \(P(w)\) is
realized among the \(n\) vertices of \(B\), and
has size at most \(st\).

Call the \(i\)-th coordinate of a profile
\emph{partial} if its size is strictly between
\(0\) and \(t\), and let \(p\) count these coordinates.
If \(2\le p<s\), its orbit has size at least
\(\binom sp t^p\ge st^2>st\).
If \(p=s\ge2\), its orbit has size at least
\(t^s>st\), with the sole exception \(s=t=2\),
where equality can occur.
For \(p=1\), if its size is \(k\),
the orbit has size
\[
       s\binom tk\binom{s-1}{r},
\]
where \(r\) is the number of the remaining
coordinates that are full. The bound \(\le st\)
forces \(k=1\) or \(k=t-1\) (so \(\binom tk=t\)),
and \(\binom{s-1}{r}=1\). Thus the other
coordinates are either all empty or all full.
The orbit consists of exactly \(st\)
distinct profiles, naturally indexed by the
\(st\) vertices of \(G\), and carries the
natural \(W\)-permutation action.

Suppose first that a partial coordinate occurs.
The \(st\) different profiles exhaust \(B\),
so the action on \(B\) is precisely the natural
action of \(S_t\wr S_s\), apart from the
exception \(s=t=2\) noted below.
There are just two orbits of \(W\) on unordered
pairs of distinct points: pairs from the same
\(t\)-block and pairs from different blocks.
Thus \(H[B]\) is either empty, complete,
\(sK_t\), or \(\overline{sK_t}\).
The edge-count identity (*) singles out \(sK_t\).

In the exceptional case \(s=t=2\), a profile
with two partial coordinates consists of
one vertex selected from each clique of \(G\).
Its orbit has four elements, with the natural
\(S_2\wr S_2\)-action on the four corners of
a square. Its two pair-orbits have sizes
four (square edges) and two (opposite corners).
By (*) the graph \(H[B]\) has exactly two edges,
so it must consist of the opposite-corner
matching \(2K_2\).

It remains to suppose \emph{every} \(P(w)\) is a
union of whole \(G\)-cliques. Then the vertices
of each \(C_i\) have the same closed neighbourhood
in \(H\): they are true twins. Because \(H\)
is vertex-transitive, all true-twin classes
have a common size \(a\ge t\).
They are cliques, and different \(C_i\)'s
lie in different classes.
Let \(Q\) be the quotient graph on these classes.
It is vertex-transitive and has
\[
                      M=\frac{2st}{a}\le2s
\]
vertices. The \(s\) classes meeting \(G\)
form an independent set of \(Q\).
If \(Q\) had no edges, \(H\) would be a
disjoint union of equal-sized cliques and
hence homogeneous, contrary to hypothesis.
Thus \(Q\) has positive degree, and its
independent sets have size at most \(M/2\).
It follows that \(s\le M/2\), and hence
\(M=2s\) and \(a=t\).
The vertices of \(G\) occupy exactly \(s\)
whole true-twin classes; the vertices of \(B\)
occupy the remaining \(s\) classes.
They already contribute \(s\binom t2\) edges
inside \(B\), which by (*) is its entire
edge count. Thus no two distinct classes
in \(B\) are adjacent and \(H[B]\cong sK_t\).
\(\square\)

\paragraph{Consequences for the manuscript.}
This lemma is independent of imprimitivity,
and so justifies applying Lemma 3.7 to
the non-homogeneous \(H\) in Proposition 3.23.
If \(G=\overline{sK_t}\), complement
both graphs first, use the lemma for
\(\overline H\) and \(\overline G=sK_t\),
and apply Lemma 3.7 there.
This also removes the invalid use of Lemma 3.3
for complements of unions of cliques.

\paragraph{Validation boundary.}
The combinatorial orbit classification was
independently enumerated for small parameters
by \texttt{checks/section3\_profiles.py}.
This is a paper-level proof, not a Lean
formalization of the lemma. No manuscript
TeX has been modified.
\end{document}
